% ============================================================
%  FCFS Simulation: Realistic vs Stress-Test Parameters
%  Supplementary comparison note
% ============================================================

\documentclass[11pt]{article}

\usepackage[margin=0.9in]{geometry}
\usepackage{amsmath}
\usepackage{booktabs}
\usepackage{tabularx}
\usepackage{array}
\usepackage{graphicx}
\usepackage{float}
\usepackage{microtype}
\usepackage{parskip}
\usepackage[hidelinks]{hyperref}
\usepackage[table]{xcolor}

\graphicspath{{figures/}}

\sloppy

\newcommand{\mainfigure}[3]{%
\begin{figure}[H]
    \centering
    \includegraphics[width=\textwidth]{#1}
    \caption{#2}
    \label{#3}
\end{figure}%
}

\title{FCFS Appointment Simulation: A Plausible Two-Class Parameterization\\[0.4em]
\large Supplementary comparison with the stress-test baseline}
\author{Alexandre Sepulveda de Dietrich\\[0.3em]
\small Under the supervision of Prof.\ Carri Chan}
\date{May 2026}

\begin{document}
\maketitle

\begin{abstract}
\noindent
This is a supplementary note to the stress-test analysis. The stress-test
baseline ($\lambda=100$ arrivals per day, $S=32$ slots, demand-to-capacity
ratio $3{:}1$) was designed to expose how simulation metrics behave under
severe capacity pressure, not to represent a plausible clinical setting.
This note re-runs the core diagnostics under a more plausible, though still
stylized, two-class parameterization drawn from the model formulation note
(v2): $\lambda=24$ arrivals per day against $S=20$ slots ($1.2{:}1$ ratio),
a $28$-day booking horizon, and behaviorally asymmetric classes representing
an MRI-like diagnostic group and a behavioral-health follow-up group.
Thirty independent seeds are used throughout.

Reducing demand pressure and introducing class asymmetry raises aggregate
utilization to $0.934$ and overall served rate to $0.789$, substantially
above the stress-test values. The more consequential change is a
$22$-percentage-point gap in class-level served rates ($88.1\%$ vs $66.1\%$)
that is absent from the symmetric stress-test baseline. This gap is
mechanically consistent with asymmetric behavioral parameters under a pooled
FCFS booking pool: Class~1 has a higher balk tolerance, lower baseline
no-show risk, and lower cancellation hazard than Class~2, and each advantage
compounds under shared-capacity FCFS. The qualitative findings of the
stress-test note survive: no-offer remains zero, offered wait exceeds
accepted wait through the same selection argument, and no-show remains the
direct slot-fill channel.
\end{abstract}

%% ============================================================
\section{Parameter Configuration}
%% ============================================================

Table~\ref{tab:params} compares the two configurations.
The most consequential differences are the demand-to-capacity ratio,
the class asymmetry, and the behavioral thresholds. Under this parameterization,
Class~1 (MRI-like) has longer balking and no-show thresholds (21~days vs
9~days and 6~days in the stress test), consistent with higher patient tolerance
for long-wait diagnostic scheduling. Class~2 (behavioral health follow-up) is
less tolerant (threshold at 14~days) and carries a higher baseline no-show
probability below the threshold ($\xi^{\text{low}}_2 = 0.15$, vs $0$ in the
stress test). Cancellation hazards are substantially lower
($0.01$--$0.02$/day vs $0.10$/day). These parameters are taken from the model
formulation note v2 and are illustrative, not empirically calibrated.

\begin{table}[h]
\centering
\renewcommand{\arraystretch}{1.18}
\small
\begin{tabular}{lcc}
\toprule
Parameter & Stress test & Comparison \\
\midrule
Total arrivals $\lambda$ (per day) & 100 & 24 \\
\quad Class 1 $\lambda_1$ & 50 & 14 \\
\quad Class 2 $\lambda_2$ & 50 & 10 \\
Slots per day $S$ & 32 & 20 \\
Demand/supply ratio & 3.1 & 1.2 \\
Booking horizon $H$ (days) & 14 & 28 \\
Burn-in / cooldown (days) & 30 / 10 & 60 / 30 \\
\midrule
\multicolumn{3}{l}{\textit{Class 1 (MRI-like diagnostic)}} \\
\quad Balk threshold (days) & 9 & 21 \\
\quad Low-delay balking prob.\ $b^{\text{low}}_1$ & 0.00 & 0.02 \\
\quad High-delay balking prob.\ $b^{\text{high}}_1$ & 0.50 & 0.55 \\
\quad No-show threshold (days) & 6 & 21 \\
\quad Low-delay no-show $\xi^{\text{low}}_1$ & 0.00 & 0.01 \\
\quad High-delay no-show $\xi^{\text{high}}_1$ & 0.30 & 0.31 \\
\quad Daily cancel prob.\ $\phi_1$ & 0.10 & 0.01 \\
\midrule
\multicolumn{3}{l}{\textit{Class 2 (Behavioral health follow-up)}} \\
\quad Balk threshold (days) & 9 & 14 \\
\quad Low-delay balking prob.\ $b^{\text{low}}_2$ & 0.00 & 0.04 \\
\quad High-delay balking prob.\ $b^{\text{high}}_2$ & 0.50 & 0.72 \\
\quad No-show threshold (days) & 6 & 14 \\
\quad Low-delay no-show $\xi^{\text{low}}_2$ & 0.00 & 0.15 \\
\quad High-delay no-show $\xi^{\text{high}}_2$ & 0.30 & 0.51 \\
\quad Daily cancel prob.\ $\phi_2$ & 0.10 & 0.02 \\
\bottomrule
\end{tabular}
\caption{Configuration comparison. Stress-test parameters are from
\texttt{configs/baseline.yaml}; comparison parameters from
\texttt{configs/realistic.yaml}, following the illustrative two-class profile
in the model formulation note v2. Both configurations are stylized.}
\label{tab:params}
\end{table}

%% ============================================================
\section{Outcome Comparison}
%% ============================================================

Table~\ref{tab:outcomes} reports aggregate and class-level outcomes for both
configurations, each averaged over 30 independent seeds.

\begin{table}[h]
\centering
\renewcommand{\arraystretch}{1.18}
\small
\begin{tabular}{lcc}
\toprule
Metric & Stress test & Comparison \\
\midrule
\multicolumn{3}{l}{\textit{Aggregate outcomes}} \\
Average utilization & 0.8418 & 0.9341 \\
Overall served rate & 0.2692 & 0.7890 \\
Mean offered wait (days) & 9.30 & 10.87 \\
Mean accepted wait (days) & 8.35 & 10.85 \\
Offered minus accepted wait (days) & 0.95 & 0.02 \\
\midrule
\multicolumn{3}{l}{\textit{Arrival outcome shares}} \\
Served & 0.269 & 0.789 \\
Balked & 0.356 & 0.033 \\
Canceled & 0.323 & 0.123 \\
No-show & 0.052 & 0.056 \\
No-offer & 0.000 & 0.000 \\
Unresolved & $<$0.001 & 0.000 \\
\midrule
\multicolumn{3}{l}{\textit{Class-level served rates}} \\
Class 1 served rate & 0.269 & 0.881 \\
Class 2 served rate & 0.269 & 0.661 \\
Class gap (C1 $-$ C2) & 0.001 & 0.220 \\
\bottomrule
\end{tabular}
\caption{Aggregate and class-level outcomes under both configurations,
averaged over 30 independent seeds. Stress-test figures are from the
controlled 30-seed diagnostic run reported in the first note.}
\label{tab:outcomes}
\end{table}

%% ============================================================
\section{What Changes and What Does Not}
%% ============================================================

\paragraph{Aggregate utilization and served rate.}
Both improve substantially relative to the stress test. Utilization rises
from $0.842$ to $0.934$; served rate from $0.269$ to $0.789$. The primary
driver of the utilization gain is the much lower cancellation hazard: with
$\phi = 0.01$--$0.02$/day instead of $0.10$/day, far fewer booked slots are
vacated before service. The approximate survival probability to service is
$0.99^{10.9} \approx 0.90$ for Class~1 and $0.98^{10.9} \approx 0.80$ for
Class~2 at mean $\tau \approx 10.9$~days, against $0.9^{8.35} \approx 0.42$
in the stress test. The served-rate gain also reflects the lower demand
pressure: the structural ceiling is $20/24 = 0.833$, and the simulation
output ($0.789$) is mechanically consistent with this bound net of no-shows.
In the stress test the ceiling is $0.32$, which the output never approaches.

\paragraph{Balking collapses; the offered/accepted gap nearly vanishes.}
In the stress test, mean offered wait ($9.30$~days) sits above the common
9-day balk threshold, driving a $35.6\%$ balked share and a $0.95$-day
offered/accepted gap through high-$\tau$ patient selection. Under this
parameterization, mean offered wait ($10.87$~days) remains well below the
Class~1 threshold (21~days) and modestly below the Class~2 threshold
(14~days). Only $3.3\%$ of arrivals balk, and the offered/accepted gap
nearly closes ($\Delta = 0.02$~days). The selection artifact that makes
accepted wait unreliable in the stress test is mechanically absent here.

\paragraph{No-show share is similar; the mechanism differs.}
No-show shares are similar ($5.2\%$ vs $5.6\%$ of arrivals), but for
structurally different reasons. In the stress test, mean accepted delay
($8.35$~days) exceeds the 6-day no-show threshold, so the fraction of
booked patients with $\tau > 6$ is substantial and those patients face
$\xi^{\text{high}} = 0.30$ no-show probability at service time. Under this
parameterization, the 21-day Class~1 threshold means that the large majority
of Class~1 bookings are placed at $\tau < 21$~days and face only
$\xi^{\text{low}}_1 = 0.01$. The aggregate no-show share is pulled up by
Class~2, for which $\xi^{\text{low}}_2 = 0.15$ applies even to short-wait
bookings. This makes Class~2 a persistent source of wasted booked capacity
across the delay distribution, not only at long delays.

\paragraph{A large class access gap emerges.}
The simulation output shows a $22$-percentage-point gap in class-level served
rates ($88.1\%$ vs $66.1\%$), which is absent from the symmetric stress-test
baseline (gap $= 0.001$) and substantially larger than the gaps produced by
the stress-test sensitivity screen under one-at-a-time parameter perturbations.
This gap is mechanically consistent with three compounding behavioral
asymmetries. Class~1 has a higher balk tolerance (threshold 21 vs 14~days),
allowing it to claim later-horizon slots that Class~2 declines. Class~1 has
a lower baseline no-show probability ($\xi^{\text{low}}_1 = 0.01$ vs
$\xi^{\text{low}}_2 = 0.15$), so its bookings are more likely to reach service.
Class~1 also has a lower daily cancellation rate ($\phi_1 = 0.01$ vs
$\phi_2 = 0.02$), reducing cumulative cancellation exposure. Under pooled FCFS,
none of these differences confer explicit priority, but all three translate
into access advantages because they affect slot consumption and booking
survival at each stage of the pipeline.

Figure~\ref{fig:class-gap} shows the class-level served rates alongside the
stress-test baseline for reference. The aggregate values look healthy; the
class gap is not visible from aggregate metrics alone.

\mainfigure{class_served_rate_gap.png}{%
Class-level served rates under the plausible two-class parameterization
(dark bars), compared with the stress-test baseline (light bars). Aggregate
access improves relative to the stress test, but access is uneven across
classes. Under pooled FCFS, this gap is mechanically consistent with
asymmetric behavioral parameters rather than explicit class priority. The
dotted line marks the structural served-rate ceiling for the comparison
parameterization ($20/24 \approx 0.833$).
}{fig:class-gap}

\paragraph{No-offer remains zero; horizon dynamics carry over.}
As in the stress test, cancellations continuously free future slots and the
horizon does not saturate, so no-offer is zero. The structural argument
is the same: daily arrivals ($24$) do not consistently outrun the combined
effect of cancellations freeing future capacity and new-slot rollover at the
horizon boundary.

%% ============================================================
\section{Summary}
%% ============================================================

The stress-test note showed that low served rate can be mechanically
consistent with high utilization when demand strongly exceeds capacity.
This comparison shows that reducing demand pressure raises aggregate served
rate and changes the loss-channel mix; balking collapses, cancellation falls,
and no-offer remains zero; but it does not eliminate class-level access
differences.

Under the current pooled FCFS formulation, the simulation output is consistent
with a $22$-percentage-point class access gap (Class~1: $88.1\%$, Class~2:
$66.1\%$) driven entirely by asymmetric behavioral parameters. Each asymmetry
; balk tolerance, baseline no-show risk, cancellation rate; operates through
the shared booking pool and compounds without any explicit class prioritization.
This gap is invisible in the aggregate utilization ($0.934$) and overall served
rate ($0.789$), which look healthy. Aggregate metrics are not sufficient to
diagnose access equity in this formulation.

The qualitative findings of the stress-test note carry over: offered wait
remains the less biased delay measure, no-show is the direct slot-fill channel,
and no-offer is zero under this parameterization. What this comparison adds is
the demonstration that FCFS neutrality does not imply outcome equality when
patient classes differ in behavioral risk, and that this effect can be
quantitatively substantial even at moderate demand-to-capacity ratios.

\end{document}
